\section{The universal determinant maps onto the matrix coefficient ring}

Fix an original setup $T$ of rank $n=2$ or $3$. Write $O$ for its
complete coefficient DVR, $k$ for its finite residue field of odd
characteristic $\ell$, $U=T.B.B$ for the universal ordinary determinant
ring, and $R=T.S.R$ for the universal normalized matrix ring. Write
$(r,\psi)$ for its matrices and ordered characters. The subring
$B=T.S.B\subset R$ is the closure, in the maximal-ideal topology of
$R$, of the $O$-algebra generated by weight-zero products of matrix
entries and all values of $\psi$. Our aim is to identify this particular
subring as the image of the universal determinant map. We keep the
specified twist and the order of the local characters throughout.

% BEGIN V68M1
\paragraph{V68M1. The twisted matrices give the required ordinary pair.}
Let $K=T.K$ be the link, with global surjection $\pi:G_0\to F.G_0$,
place bijection $e$, local group isomorphisms $\lambda_v$, and scalar
character $\kappa:G_0\to O^\times$. Write $\phi_v$ for the original
local-to-global map and $J$ for the antidiagonal permutation matrix.
For the universal lift put
\[
 q(g)=\kappa(g)r(\pi g),\qquad
 b_{vi}(s)=\psi_{e(v)i}(\lambda_v s)\kappa(\phi_v s).
\]
Scalars in $O$ act through the given coefficient map. The ordinary
pair over $R$ is the degree-$n$ determinant of $q$ and the ordered
characters $b$. To verify all its conditions, extend $q$ linearly to
the group algebra, writing $q(y)$ for the resulting matrix. After
any scalar extension the characteristic polynomial of $y$ for this
determinant is $\det(XI-q(y))$, by evaluating the defining polynomial
law at $X[1]-y$. Its first signed coefficient is $\operatorname{tr}q(y)$.
Modulo each power of the maximal ideal, the finitely many entries of
$q(g)$ are simultaneously locally constant by the lift's continuity,
continuity of $\pi$, and adic continuity of $\kappa$. All signed
characteristic coefficients are therefore adically continuous. The
same product argument gives continuity of $b_{vi}$ on the local group.
Reduction of $q$ is the diagonal matrix of the prescribed characters
$\bar\chi_i$, since the lift reduces diagonally and the link identifies
$\bar\kappa\,\Psi^{w_i}\circ\pi$ with $\bar\chi_i$. Reducing the
polynomial determinant identity consequently gives the prescribed
residual law. Likewise $b_{vi}$ reduces to $\bar\chi_i\circ\phi_v$.
On the specified subgroup $\mu_v$ the link places every element in
the closure of the commutator subgroup. An adically continuous
character is one there, since the fiber of one is closed modulo each
maximal-ideal power and their intersection is zero. The prescribed
character $\theta_{vi}$ is also one there by its defining condition.

Write $c(g)$ for conjugation from the determinant frame. Its group
algebra star takes $[g]$ to $\epsilon(g)^{1-n}[c(g)^{-1}]$. Since
$\kappa(g)\kappa(c(g))=\epsilon(g)^{1-n}$ and $\pi$ respects
conjugation, matrix polarization gives $q([g]^*)=Jq(g)^tJ$. This
preserves every characteristic coefficient. For any ordered tuple
$s_i$ in the local group, the matrix image of the ordinary relation is
\[
 \prod_{i=0}^{n-1}\bigl(q(\phi_v s_i)-b_{vi}(s_i)I\bigr)=
 \Bigl(\prod_i\kappa(\phi_v s_i)\Bigr)
 \prod_{i=0}^{n-1}\bigl(r(\pi\phi_v s_i)-
                         \psi_{e(v)i}(\lambda_v s_i)I\bigr)=0.
\]
The link respects the local maps, so the ordered product on the right
vanishes by the normalized ordinary relation. Finally take $g$ in an
auxiliary inertia group with every $\bar\chi_i(g)=1$. The tameness condition
of the link gives $m$ prime to $\ell$ with $(\pi g)^m=1$ and
$\kappa(g)^m=1$. Thus $q(g)^m=I$ and $q(g)\equiv I$ in the residue
field. The geometric sum $I+q(g)+\cdots+q(g)^{m-1}$ is invertible,
since its reduction is $mI$. Multiplying its identity with $q(g)-I$
gives $q(g)=I$. Both required trace relations now follow by left
multiplying the vanishing matrices by any $q(h)$ and taking traces.
This verifies the original $\mathrm{Pair}$, including its auxiliary
inertia, local torsion, polarization and ordered ordinary conditions.
% END V68M1

% BEGIN V68M2
\paragraph{V68M2. Descent of the whole determinant law to $B$.}
For a finite list $g_1,\ldots,g_t$ in $G_0$, form
\[
 P_{\boldsymbol g}(X_1,\ldots,X_t)
        =\det\Bigl(\sum_{a=1}^t X_a q(g_a)\Bigr).
\]
All its coefficients belong to $B$. Indeed a determinant term is a
product of $n$ scalars from $O$ and entries
$r(\pi g_{a_i})_{i,\sigma(i)}$ indexed by a permutation $\sigma$.
The entry weights sum to $\sum_i(w_i-w_{\sigma(i)})=0$.
Every such product is in the defining algebra of $B$, and signed sums
remain there. In particular this includes all coefficients of
$\det(XI-q(g))$ by using $1,g$ and the substitution $(X,-1)$.

For any commutative $B$-algebra $H$ define a determinant on $H[G_0]$
by evaluating $P_{\boldsymbol g}$ at the coefficients of a finite
group-algebra sum. This definition is independent of the list by the
polynomial identities for permutation, combining repeated group
elements and adjoining a zero coefficient. They hold over $R$ and
hence over $B$, by injectivity on polynomial coefficients. The same
argument for two lists proves multiplicativity: use independent
variables $X_a,Y_b$ and the product list $g_a h_b$ with coefficients
$X_aY_b$. Over $R$ this is the determinant of a product of matrices.
The identity therefore holds in $B[\boldsymbol X,\boldsymbol Y]$ and
can be specialized in any $H$. The law sends $[1]$ to one, is natural
in $H$, and is homogeneous of degree $n$. These identities define the
full group determinant over $B$, with scalar extension to $R$ equal to
the determinant of $q$ on all group-algebra elements.

The values of $b_{vi}$ and their inverses belong to $B$ by the
definition of $B$ and evaluation at $s^{-1}$. By RTC6--RTC7, $B$
is a complete Noetherian local coefficient ring with the same residue
field as $R$. For each $N$, RTC8 in degree zero gives an $M$ with
$B\cap\mathfrak m_R^M\subset\mathfrak m_B^N$. Thus V68M1 gives
adic continuity over $B$ for both the signed coefficients and the
characters. The residue maps of $R$ and $B$ agree on $B$, so reducing
the generic polynomials gives the prescribed residual law over every
$k$-algebra. All other pair identities in V68M1 descend through
$B\hookrightarrow R$. The universal property of $U$ applied to this
pair supplies an $O$-algebra map $\Phi:U\to B$. After inclusion in $R$
it carries the universal determinant law and the ordered characters
to those of V68M1. This is the exact condition $T.\mathrm{IsPhi}(\Phi)$.
% END V68M2

% BEGIN V68M3
\paragraph{V68M3. Recovering weight-zero products from traces in ranks two and three.}
Let $C$ be the image in $R$ of a map from a complete Noetherian local
$O$-coefficient algebra $C_0$ with residue field $k$. Suppose $C$
contains every $\operatorname{tr} r(g)$ and every local character
value $\psi_{vi}(s)$. We prove $B\subset C$. The quotient
$C_0/\ker(C_0\to R)$ makes $C$ a complete Noetherian local coefficient
algebra with the same residue field. Indeed the map is local by RTC1,
and quotients by ideals are complete in this setting as in RTC6.
Moreover $C$ is closed in $R$: the source is compact as the inverse
limit of its finite adic quotients (each successive layer is finite
dimensional over the finite field $k$), the map is continuous by RTC1,
and $R$ is Hausdorff.

Let $V$ be the $C$-span in $M_n(R)$ of all $r(g)$. It contains $I$,
is closed under products by the group law, and every element has trace
in $C$. At the marked element $g_0$ the difference
$r(g_0)-r(g_0^*)$ is diagonal by normalization. Polarization reverses
its diagonal entries and changes their sign when taking this
difference. Hence the difference is $zH$, where $H=\operatorname{diag}(w_i)$
and the weights are $(1,-1)$ or $(1,0,-1)$. Its first entry $z$ reduces
to $\Psi(g_0)-\Psi(g_0)^{-1}\ne0$ by the marked-frame condition.
Since $zH\in V$, its square has trace $2z^2\in C$. The residue
characteristic is odd, so $z^2\in C$ and $2$ is a unit there.
Hensel's lemma in $C$ lifts the residue of $z$ to a square root $d$
of $z^2$: the derivative at the residue root is a unit. In $R$,
$(d-z)(d+z)=0$ and $d+z$ has nonzero residue, so $d=z$. Consequently
$z\in C^\times$ and $H\in V$. The projections to weights $1$ and
$-1$ are $(H^2+H)/2$ and $(H^2-H)/2$. In rank three the remaining
projection is $I-H^2$. Thus every diagonal matrix unit $E_{ii}$ is in $V$.

Write $V_{ij}$ for the $C$-span of entries $r(g)_{ij}$ in $R$. For
$x\in V_{ij}$ the matrix $xE_{ij}$ is in $V$ by left and right
multiplication by the projections. Hence $V_{ii}\subset C$ by trace,
and $V_{ij}V_{jt}\subset V_{it}$ by matrix multiplication followed by
taking the $(i,t)$ entry. Polarization exchanges $V_{ij}$ and
$V_{j^*i^*}$, where reversal $i\mapsto i^*$ changes the sign of the
weight. A product of total weight zero now reduces by induction on
its length. A zero-weight factor is diagonal and lies in $C$. If all
remaining weights are nonzero, select factors of opposite sign. In
rank two their endpoints are reversed, so they compose. In rank three
an edge of degree $1$ can be represented as either $(1,0)$ or
$(0,-1)$ by reversal, and one of degree $-1$ as either $(0,1)$ or
$(-1,0)$. The degree $2$ and $-2$ edges are $(1,-1)$ and $(-1,1)$.
Thus two opposite-sign edges can be made composable, in one of the
two orders, also when one or both have degree two in absolute value.
Their product is a single entry-space factor of the sum degree.
The induction begins with the empty product. Every balanced monomial
therefore belongs to $C$. Including the local character values and
using closedness in $R$ proves $B\subset C$ with the original adic
closure in its definition.
% END V68M3

% BEGIN V68M4
\paragraph{V68M4. Surjectivity on the original coefficient subring.}
Apply V68M3 to the image of $U\xrightarrow{\Phi} B\hookrightarrow R$.
For any $x\in F.G_0$, choose $g$ with $\pi g=x$. The universal first
signed coefficient at $[g]$ maps to $\operatorname{tr}q(g)$ by V68M2.
Multiplication by the inverse of the $O$-unit $\kappa(g)$ therefore
puts $\operatorname{tr}r(x)$ in the image. For each original local
place and element, use the surjectivity of $e$ and of $\lambda_v$ to
write them as $e(v)$ and $\lambda_v s$. The image contains the value
$b_{vi}(s)$ of the universal ordered character and the inverse of the
scalar $\kappa(\phi_v s)$, hence also $\psi_{e(v)i}(\lambda_v s)$.
V68M3 gives $B$ inside this image, while V68M2 gave $\Phi$ with
codomain $B$. It follows that $\Phi$ is surjective. Together with
$T.\mathrm{IsPhi}(\Phi)$ this proves the full original $T.\mathrm{PhiEx}$
on each supplied rank-two or rank-three setup.
% END V68M4
